% ======================================================================
%  sec22.tex —— 第 22 节 （局部）诺特概形
%  底本：原书 pp.108–110（印刷页 = PDF 页）
%  原书标题：22. (Locally) Noetherian schemes
% ======================================================================

\section{（局部）诺特概形}\label{sec:22}

% ---------------------------------------------------------------- p.108
\subsection{定义}\label{sec:22.1}

\begin{definition}\label{def:22.1.1}
称概形 $X$ 为\emph{局部诺特的}，若对任意仿射开子概形 $\Spec(A) \subseteq X$，交换环 $A$ 都是诺特环。称 $X$ 为\emph{诺特的}，若它局部诺特且拟紧。
\end{definition}

\begin{yizhu}
「每个仿射开子概形」这一要求看似过强：概形只由一族仿射开覆盖给出，覆盖中出现的环诺特，未必能保证其余仿射开子概形上的环也诺特。真正省力的是下一引理——只要\emph{存在}一族仿射诺特覆盖就够。Stacks Project 的定义正取此形式，并把它与「每点有仿射诺特开邻域」等同。
[来源：Stacks Project, Tag 01OV（Definition 28.5.1，诺特概形），\url{https://stacks.math.columbia.edu/tag/01OV}（2026-11-17 访问）。]
\end{yizhu}

\begin{lemma}\label{lem:22.1.2}
设 $X$ 为概形。若存在由仿射开子概形给出的覆盖 $X = \bigcup_{i \in I} X_i$，使得对每个 $i \in I$，$\cO_X(X_i)$ 都是诺特环，则 $X$ 局部诺特。
\end{lemma}

\begin{proof}
设 $U \subseteq X$ 为仿射开子概形。取 $U$ 的开覆盖 $U = \bigcup_{j \in J} U_j$，使每个 $U_j$ 都是某个 $X_i$ 的标准仿射开子概形。于是 $\cO_X(U_j)$ 诺特，因为它是诺特环的局部化。因此可把 $X$ 换成 $U$，假设 $X$ 仿射。又 $X$ 拟紧，故可进一步设 $I$ 有限。

记 $A = \cO_X(X)$，$A_i = \cO_X(X_i)$。设 $J_1 \subseteq J_2 \subseteq \cdots$ 为 $A$ 的理想升链，
\[
J_1 \subseteq J_2 \subseteq \cdots \tag{22.1}\label{eq:22.1}
\]
是对应的 $\cO_X$ 的拟凝聚理想升链。只需证明此链稳定。

对任意 $i \in I$，式 (22.1) 的限制
\[
J_1|_{X_i} \subseteq J_2|_{X_i} \subseteq \cdots
\]
是 $\cO_{X_i}$ 的拟凝聚理想升链。由于 $X_i \simeq \Spec(A_i)$ 且 $A_i$ 诺特，该链必须稳定。又 $I$ 有限，故链 (22.1) 也稳定。
\end{proof}

\begin{yizhu}
nLab 的条目直接以「可被仿射诺特开子集覆盖」作为局部诺特概形的定义，并补一句：诺特环的局部化仍诺特，故这一覆盖条件自动升级为「每个仿射开子概形都如此」。这正是引理 \ref{lem:22.1.2} 的算术内核，几何部分只剩「换到标准仿射开子概形」这一步归约。
[来源：nLab "noetherian scheme"，\url{https://ncatlab.org/nlab/show/noetherian+scheme}（Last revised 2014-01-11；2026-11-17 访问）。]
\end{yizhu}

\begin{corollary}\label{cor:22.1.3}
仿射概形 $\Spec(A)$ 诺特当且仅当 $A$ 诺特。
\end{corollary}

\begin{corollary}\label{cor:22.1.4}
局部诺特性在概形上是局部的。
\end{corollary}

\begin{warning}\label{warn:22.1.5}
诺特概形的纤维积未必局部诺特。见例 \ref{ex:0.2.3}。
\end{warning}

\begin{yizhu}
纤维积 $\Spec(\overline{\Q}) \times_{\Spec(\Q)} \Spec(\overline{\Q})$ 的底概形是 $\Spec(\overline{\Q} \otimes_{\Q} \overline{\Q})$。两个因子都是诺特概形，而该张量积不是诺特环（例 \ref{ex:0.2.3}），故这个纤维积连局部诺特都不是。可见诺特性对基变换不稳定；稳定的版本要求基概形局部诺特、态射有限型，见引理 \ref{lem:23.1.7}。
[依据：原书 p.108 警告 22.1.5；张量积的反例见原书例 0.2.3（p.7）。]
\end{yizhu}

\subsection{诺特空间}\label{sec:22.2}

\begin{definition}\label{def:22.2.1}
称拓扑空间 $X$ 为\emph{诺特的}，若 $X$ 的闭子集满足降链条件，即任意闭子集降链都稳定。称 $X$ 为\emph{局部诺特的}，若对任意点 $x \in X$，$x$ 在 $X$ 中都有诺特开邻域。
\end{definition}

\begin{exercise}\label{exc:22.2.2}
考虑（局部）诺特拓扑空间构成的类。
\begin{enumerate}[label=(\roman*),leftmargin=2em,itemsep=0.25ex]
  \item 空间诺特当且仅当它局部诺特且拟紧。
  \item （局部）诺特空间的子空间（局部）诺特。
  \item 有限多个（局部）诺特子空间之并（局部）诺特。
  \item 空间诺特当且仅当它的任意开子空间都拟紧。
\end{enumerate}
\end{exercise}

\begin{lemma}\label{lem:22.2.3}
诺特拓扑空间 $X$ 只有有限多个不可约分支，且每个不可约分支都包含 $X$ 的一个非空开子集。
\end{lemma}

\begin{proof}
先证第一条。设 $\mathcal{I}$ 为 $X$ 的含无穷多个不可约分支的闭子空间构成的集合。只需证 $\mathcal{I} = \varnothing$。若 $\mathcal{I}$ 非空，由 $X$ 诺特知 $\mathcal{I}$ 有极小元 $Y$。按 $\mathcal{I}$ 的定义，$Y$ 不是不可约的。于是可写 $Y = Y_1 \cup Y_2$，其中 $Y_1, Y_2 \ne Y$。由 $Y$ 的极小性，$Y_i \notin \mathcal{I}$，故每个 $Y_i$ 都是有限多个不可约子空间之并，从而 $Y$ 也是。容易看出这意味着 $Y$ 只有有限多个不可约分支。矛盾。

再证第二条。设 $Z_1, \dots, Z_n$ 为 $X$ 的不可约分支。容易看出 $X \setminus (Z_2 \cup \cdots \cup Z_n)$ 是 $X$ 的含于 $Z_1$ 的非空开子集。
\end{proof}

\begin{yizhu}
Stacks Project 把这条引理与「诺特空间的子空间仍诺特」并列为同一条目，后者的证明只用闭子集的降链条件，属于习题 \ref{exc:22.2.2}(ii)。另注意第二条断言不能加强为「分支本身开」：$Z_1$ 含开集 $X \setminus (Z_2 \cup \cdots \cup Z_n)$，但 $Z_1$ 未必开，例如 $\Spec(k[x,y]/(xy))$ 的两个分支。
[来源：Stacks Project, Tag 0052（Lemma 5.9.2，诺特空间的不可约分支），\url{https://stacks.math.columbia.edu/tag/0052}（2026-11-17 访问）。]
\end{yizhu}

\begin{lemma}\label{lem:22.2.4}
（局部）诺特概形的底拓扑空间是（局部）诺特的。
\end{lemma}

\begin{proof}
由习题 \ref{exc:22.2.2}，只需证诺特概形的底拓扑空间局部诺特，从而只需证 $\Spec(A)$ 的底拓扑空间诺特。但这是显然的。
\end{proof}

\begin{yizhu}
反向不成立。取 $A = k \oplus V$，其中 $V$ 为无限维 $k$-向量空间，乘法规定 $V \cdot V = 0$。则 $A$ 只有唯一素理想 $V$，故 $\Spec(A)$ 是单点空间，当然诺特；但 $V$ 作为理想不有限生成，$A$ 不是诺特环。「底空间诺特」远弱于「概形诺特」。
[译者按（上下文推断，无外部依据）：这是交换代数中的标准反例，用以说明拓扑条件不能替代环条件。]
\end{yizhu}

\begin{exercise}\label{exc:22.2.5}
设 $X$ 为局部诺特概形，$x \in X$ 为一点。证明：若 $\Spec(\cO_{X,x})$ 不可约，则存在 $x$ 的开邻域 $U$，使 $U$ 不可约。
\end{exercise}

\subsection{诺特概形的浸入}\label{sec:22.3}

\begin{proposition}\label{prop:22.3.1}
设 $f\colon X \to Y$ 为概形的局部闭浸入。若 $Y$ （局部）诺特，则 $X$ 亦然。
\end{proposition}

\begin{proof}
只需在 $f$ 为闭浸入或开浸入时证明。此外，两种情形下的断言都在 $Y$ 上局部，故可设 $Y = \Spec(B)$，其中 $B$ 诺特。

闭浸入的情形显然，因为 $B$ 的任意商环都诺特。开浸入的情形：设 $f\colon X \to Y$ 为开浸入。由推论 \ref{cor:22.1.4}，$X$ 局部诺特。由引理 \ref{lem:22.2.4} 与习题 \ref{exc:22.2.2}，$X$ 的底拓扑空间拟紧。故 $X$ 是诺特概形。
\end{proof}

\begin{lemma}\label{lem:22.3.2}
设 $X$ 为局部诺特概形。任意态射 $f\colon X \to Y$ 都拟分离。
\end{lemma}

\begin{proof}
由引理 \ref{lem:18.3.4}(ii)，只需证 $X$ 拟分离。由引理 \ref{lem:18.3.3}，只需证对任意仿射开子集 $U, V \subseteq X$，交 $U \cap V$ 拟紧。由假设 $U$ 诺特。由命题 \ref{prop:22.3.1}，$U \cap V$ 作为 $U$ 的开子概形诺特，特别地拟紧。
\end{proof}

\begin{lemma}\label{lem:22.3.3}
设 $X$ 为诺特概形，$Y$ 为局部诺特概形。任意态射 $f\colon X \to Y$ 都拟紧。
\end{lemma}

\begin{proof}
由引理 \ref{lem:18.1.4}(ii) 与引理 \ref{lem:22.3.2} 即得。
\end{proof}

由命题—定义 \ref{prop:18.5.4} 与引理 \ref{lem:18.5.5}，我们得到：

% ---------------------------------------------------------------- p.110
\begin{corollary}\label{cor:22.3.4}
设 $Y$ 为局部诺特概形，$i\colon X \to Y$ 为概形的局部闭浸入。则 $X$ 在 $Y$ 中的闭包 $\bar X$ 良定义，且典范态射 $j\colon X \to \bar X$ 是开浸入。
\end{corollary}

\subsection{局部诺特概形上的凝聚模}\label{sec:22.4}

\begin{lemma}\label{lem:22.4.1}
设 $X$ 为局部诺特概形，$\cF$ 为 $\cO_X$-模。则下列条件等价：
\begin{enumerate}[label=(\roman*),leftmargin=2em,itemsep=0.25ex]
  \item $\cO_X$-模 $\cF$ 凝聚；
  \item $\cO_X$-模 $\cF$ 有限表现；
  \item $\cF$ 拟凝聚且有限型；
  \item 对任意仿射开子概形 $U \subseteq X$，有 $\cF|_U \simeq \widetilde M$，其中 $M$ 是某个有限生成 $\cO_X(U)$-模。
\end{enumerate}
\end{lemma}

\begin{proof}
(i)$\Rightarrow$(ii) 由习题 \ref{exc:16.2.5}(1) 得到。(ii)$\Rightarrow$(iii) 由引理 \ref{lem:16.2.2} 与引理 \ref{lem:16.2.3} 得到。

(iii)$\Rightarrow$(iv)：可设 $X = U = \Spec(A)$ 仿射，且 $\cF \simeq \widetilde M$，其中 $M$ 为 $A$-模。只需证 $M$ 有限生成。因 $\cF$ 有限型，可用标准仿射开子集覆盖 $X = \bigcup_{i \in I} U(f_i)$，使 $\cF|_{U(f_i)}$ 由有限多个截面生成。这意味着 $M_{f_i}$ 是有限生成 $A_{f_i}$-模。由习题 \ref{exc:16.1.8}，$M$ 有限生成。

(iv)$\Rightarrow$(i)：为验证定义 \ref{def:16.4.1} 中的条件，只需处理属于 $X$ 的某个拓扑基的开子集，故可设 $U$ 仿射。此时条件由 (iv) 立即得到，因为 $\cO_X(U)$ 诺特。
\end{proof}

\begin{yizhu}
在一般概形上「凝聚」严格强于「有限表现」：后者只要求局部有有限自由表现，前者还要求任意开集上的核条件成立。Stacks Project 的对应引理正是在局部诺特假设下把这四条打通，并顺带得出结构层 $\cO_X$ 自身凝聚，这是 $\cO_X(U)$ 诺特的直接推论。
[来源：Stacks Project, Tag 01XZ（Lemma 30.9.1，局部诺特概形上的凝聚模），\url{https://stacks.math.columbia.edu/tag/01XZ}（2026-11-17 访问）。]
\end{yizhu}

\begin{corollary}\label{cor:22.4.2}
「凝聚模」这一性质在局部诺特概形上是局部的。
\end{corollary}

\begin{proof}
引理 \ref{lem:22.4.1} 的条件 (iii) 在概形上是局部的。
\end{proof}

\begin{corollary}\label{cor:22.4.3}
设 $f\colon X \to Y$ 为局部诺特概形之间的态射。则 $f^*$ 保持凝聚模。
\end{corollary}

\begin{corollary}\label{cor:22.4.4}
设 $X$ 为局部诺特概形。$\cO_X$-凝聚模对取核、余核与扩张封闭。特别地，它们构成阿贝尔范畴 $\mathsf{Coh}(X)$，且嵌入函子 $\mathsf{Coh}(X) \subseteq \mathsf{QCoh}(X)$ 正合。
\end{corollary}

\begin{yizhu}
Stacks Project 把这三条封闭性合并为「凝聚模范畴是阿贝尔范畴」，并指出关键仍在局部诺特：此时凝聚性与有限表现性一致，而有限表现模对核、余核、扩张封闭是形式事实。正合性则来自 $\mathsf{Coh}(X) \subseteq \mathsf{QCoh}(X)$ 作为全子范畴的嵌入恰好保持这两类极限。
[来源：Stacks Project, Tag 01Y0（Lemma 30.9.2，局部诺特概形上的凝聚模），\url{https://stacks.math.columbia.edu/tag/01Y0}（2026-11-17 访问）。]
\end{yizhu}

\begin{remark}\label{rem:22.4.5}
若 $X$ 诺特，则 $\mathsf{Coh}(X) \subseteq \mathsf{QCoh}(X)$ 是紧对象构成的全子范畴。
\end{remark}

\begin{yizhu}
「紧对象」是范畴论用语：对象 $C$ 紧，指函子 $\Hom(C, -)$ 与滤过余极限交换。nLab 提醒，这一范畴意义与拓扑紧性并不同义——带 indiscrete 拓扑的两点空间是紧空间，却不是 $\mathsf{Top}$ 的紧对象。把 $\mathsf{Coh}(X)$ 认作 $\mathsf{QCoh}(X)$ 的紧对象，正是「有限表现即与滤过余极限交换」的几何版本。
[来源：nLab "compact object"，\url{https://ncatlab.org/nlab/show/compact+object}（Last revised 2026-05-22；2026-11-17 访问）。]
\end{yizhu}

\begin{exercise}\label{exc:22.4.6}
设 $X$ 为局部诺特概形，$M, N \in \mathsf{Coh}(X)$。则 $M \otimes_{\cO_X} N \in \mathsf{Coh}(X)$。
\end{exercise}
